\documentclass[10pt,a4paper]{amsart}
\usepackage[margin=19mm]{geometry}
\usepackage{amsmath,amssymb,mathtools}
\usepackage[scr=rsfs,cal=euler]{mathalfa}
\usepackage{xfrac}
\usepackage[unicode,hidelinks,bookmarks=false]{hyperref}
\newcommand{\ie}{i.e.\ }
\newcommand{\cf}{cf.\ }
\newcommand{\resp}{respectively\ }
\newcommand{\Subsection}{\subsection}
\providecommand{\nobreakdash}{\nobreak\hbox{-}}
\setlength{\emergencystretch}{2em}
\multlinegap0pt
\usepackage[shortlabels]{enumitem}
\usepackage{amssymb}
\usepackage{cancel}
\usepackage{tikz}
\newtheorem{lemma}{Lemma}
\usepackage{cleveref}
\crefname{lemma}{Lemma}{Lemmas}
\newcommand{\PG}{\mathrm{PG}}
\newcommand{\fq}{{\mathbb F}_{q}}
\newcommand{\fqn}{\mathbb{F}_{q^n}}
\renewcommand{\mod}{\hbox{{\rm mod}\,}}
\newcommand{\pg}{\PG}
\DeclareMathOperator{\ww}{w}
\title{On the minimum size of linear sets}
\author{Sam Adriaensen, Paolo Santonastaso}
\date{Source: arXiv:2301.13001v2 (CC BY 4.0)}
\begin{document}
\maketitle

\noindent\emph{Source and licence.} On the minimum size of linear sets. Version arXiv:2301.13001v2 (CC BY 4.0), distributed by arXiv under the Creative Commons Attribution 4.0 International licence, http://creativecommons.org/licenses/by/4.0/, as recorded in the arXiv licence metadata for this exact version and shown on the abstract page https://arxiv.org/abs/2301.13001v2. Full licence notice: \texttt{licenses/arxiv-2301.13001-CC-BY-4.0.txt}.\par\smallskip

\noindent\textit{Editorial scope.} Counted unit: the article's lemma lm:subgeometry (source0.tex lines 446--450). The article's complete original proof of it is delivered with it (source0.tex lines 452--493, one \texttt{proof} environment of 42 lines). No mathematics is written, repaired or re-derived: the statement and the proof are the article's own lines, in the article's own notation, and no retained line is substituted or re-flowed.  Retained uncounted as the source-local context the proof needs: lines 306--311; lines 355--363.  Nothing was omitted from any retained span, so the package carries no omission record.  No retained line contains a citation, so the wrapper carries no bibliography and no bibliography entry.  No retained line contains a figure environment, an image-inclusion command or drawing code, so no image is delivered and the package declares no asset dependency.  Editorial formatting only, in addition: an AMS \texttt{amsart} preamble that restates the article's own declarations and macros used by the retained lines, namely the environments \texttt{lemma} on separate counters, together with the standard packages those lines need.  The retained environments print in this document as Lemma 1 in order of appearance, whereas the article prints the same items with the numbering of its own full document.  The complete native source of the article as distributed in the arXiv e-print for this exact version is bundled unchanged as \texttt{sources/136703/source0.tex}.
\begin{align}
    |L_U| &=N_1+\ldots+N_n, \label{eq:pesicard} \\
    \sum_{i=1}^n N_i \frac{q^i-1}{q-1}&=\frac{q^k-1}{q-1}, \label{eq:pesivett} \\
    |L_U| &\leq \frac{q^k-1}{q-1}, \label{eq:card} \\
    |L_U| &\equiv 1 \ (\mod q). \label{eq:modqsize}
\end{align}

\begin{lemma}
 \label{lm:IndependentSubspaces}
 Let $W_1, \dots, W_m$ be independent subspaces in $\fqn^{d+1}$, and let $L_U$ be an $\fq$-linear set in $\PG(d,q^n)$ of rank $k$, that spans the entire space.
 Then
 \[
  \ww_{L_U}( \pg(W_1,\fqn)) + \ldots + \ww_{L_U}(\pg(W_m,\fqn)) \leq k.
 \]
 If equality holds, then $\fqn^{d+1} = W_1 \oplus \ldots \oplus W_m$.
\end{lemma}

\begin{lemma}
 \label{lm:subgeometry}
Let $L_U$ be an $\fq$-linear set in $\PG(d-1,q^n)$, with $d \geq 2$.
Then $L_U$ spans $\PG(d-1,q^n)$ and satisfies $|L_U| = \frac{q^d-1}{q-1}$ if and only if $L_U$ is a canonical $\fq$-subgeometry.
\end{lemma}

\begin{proof}
If $L_U$ is a canonical subgeometry, then it immediately follows that $L_U$ spans the entire space, and $|L_U| = \frac{q^d-1}{q-1}$.
So suppose that $L_U$ spans the space, and that $|L_U| = \frac{q^d-1}{q-1}$.
We need to prove that all points of $L_U$ have weight 1.
Indeed in that case, by equations \eqref{eq:pesicard} and \eqref{eq:pesivett}, $L_U$ must then have rank $d$, which proves that $L_U$ is a canonical subgeometry.
So suppose by way of contradiction that $L_U$ has points of weight greater than 1.
Note that by Equations \eqref{eq:pesicard} and \eqref{eq:pesivett}, the rank of $L_U$ is some number $k > d$.
Let
\[
 \sigma = \langle P \in L_U \colon \ww_{L_U}(P) > 1 \rangle
\]
denote the subspace of $\PG(d-1,q^n)$ spanned by points of weight greater than 1.

Suppose that $\sigma$ is not $\PG(d-1,q^n)$. 
Then $L_U \not \subseteq \sigma$, and every point in $L_U \setminus \sigma$ is a point of weight 1.
Hence, there are at least $q^{k-1}$ points in $L_U \setminus \sigma$ corresponding to (necessarily distinct) points of weight 1 of $L_U$.
Thus, $|L_U| > q^{k-1} > \frac{q^d-1}{q-1}$ since $k > d$, a contradiction.

Hence $\sigma$ equals $\PG(d-1,q^n)$.
Let
\[
 m = \max_{P \in L_U} \ww_{L_U}(P)
\]
denote the maximum weight of the points of $L_U$.
Then we can choose $d$ independent points $P_1, \dots, P_d$ in $L_U$ such that $\ww_{L_U}(P_1) = m$, and $\ww_{L_U}(P_i) \geq 2$ for each $i$.
By \Cref{lm:IndependentSubspaces},
\[
 k \geq \sum_{i=1}^d \ww_{L_U}(P_i) \geq m + 2(d-1).
\]
Let $N_1,\ldots,N_m$ denote the weight distribution of $L_U$.
Then by Equations \eqref{eq:pesicard} and \eqref{eq:pesivett},
\[
 \frac{q^m-1}{q-1} |L_U| = \sum_{i=1}^m N_i \frac{q^m-1}{q-1}
 \geq \sum_{i=1}^m N_i \frac{q^i-1}{q-1} = \frac{q^k-1}{q-1}
 \geq \frac{q^{m+2(d-1)}-1}{q-1}.
\]
This implies that
\[
  \frac{q^d-1}{q-1} = |L_U| \geq \frac{q^{m+2(d-1)}-1}{q^m-1},
\]
which yields a contradiction if $d \geq 2$.
\end{proof}

\end{document}
